\conclusion
В работе систематизированы теоретические и прикладные аспекты МКВ. Криптографические основы обеспечения корректности и приватности изложены в разделе~\ref{sec:cryptographic-foundations}, классификация по числу участников и моделям безопасности --- в разделе~\ref{sec:protocol-classification}, применения, включая анализ зарплат, обучение нейронных сетей и управление криптовалютными ключами, --- в разделе~\ref{sec:applications}.

Сопоставление семейств протоколов приведено в таблице~\ref{tab:mpc-primitives}, ограничения вычислений, сетевого обмена и обработки дробных чисел --- в разделе~\ref{sec:limitations}. Работа Zhang et al. показывает сокращение числа умножений при сохранении точности за счёт оптимизации свёртки и Softmax на основе схемы Шамира (см. подраздел~\ref{subsec:ml-applications}); эффективность нелинейных функций остаётся нерешённой задачей.

Для будущего исследования предложен продукт из двух раздельных блоков: логистической регрессии и заменяемого МКВ-модуля. Выбраны TFEncrypted и ABY3 для трёх серверов в semi-honest модели при честном большинстве. Приоритет --- простота реализации, а не максимальная эффективность или защищённость (см. раздел~\ref{sec:future-research}).

Цель MVP --- обучение и инференс без раскрытия входов с проверкой численных допусков относительно открытой контрольной модели (см. подраздел~\ref{subsec:mvp-implementation-plan}). Реализация и измерения предстоят. Демонстрация не заменит доказательство безопасности; защита от активного противника, fairness и промышленная готовность не заявляются.

После MVP предусматриваются расширение испытаний и отдельное исследование DP, TEE, FL и ORAM (см. подраздел~\ref{subsec:mvp-risks-roadmap}). Замена МКВ-модуля при сохранении границы блоков потребует повторной проверки численной совместимости и модели угроз.